% ============================================================
%  Informe de Arquitectura Multiagente para Teoría del Caso Aumentada
%  Formato APA 7ª edición
%  Autores: Santiago Esteban Redondo Perdomo, Juan David Conde Duarte
%  Materia: Ciencia de Datos
% ============================================================
\documentclass[12pt, letterpaper]{article}

% ---- Paquetes esenciales ----
\usepackage[utf8]{inputenc}
\usepackage[T1]{fontenc}
\usepackage{times}            % Times New Roman (aproximación APA)
\usepackage[margin=1in]{geometry}
\usepackage{setspace}
\usepackage{parskip}          % Espacio entre párrafos, sin sangría
\usepackage{hyperref}
\usepackage{enumitem}
\usepackage{fancyhdr}
\usepackage{graphicx}
\usepackage{booktabs}
\usepackage{tabularx}
\usepackage{xcolor}
\usepackage{listings}
\usepackage{titlesec}

% ---- Formato APA ----
\doublespacing
\setlength{\parindent}{0.5in}

% ---- Encabezado (running head) ----
\pagestyle{fancy}
\fancyhf{}
\fancyhead[L]{\small\textsc{Arquitectura Multiagente para Teoría del Caso}}
\fancyhead[R]{\thepage}
\renewcommand{\headrulewidth}{0pt}

% ---- Portada ----
\title{\textbf{Informe de Arquitectura Multiagente para Teoría del Caso Aumentada}}
\author{
  Santiago Esteban Redondo Perdomo \\
  Juan David Conde Duarte \\
  \\
  \textit{Ciencia de Datos}
}
\date{}

% ---- Formato de títulos APA (sin negrita excesiva) ----
\titleformat{\section}{\normalfont\bfseries\centering}{}{0em}{}
\titleformat{\subsection}{\normalfont\bfseries\centering}{}{0em}{}
\titleformat{\subsubsection}{\normalfont\bfseries}{}{0em}{}

% ---- Configuración de listings ----
\lstset{
  basicstyle=\footnotesize\ttfamily,
  breaklines=true,
  frame=single,
  numbers=none,
  captionpos=t,
  aboveskip=12pt,
  belowskip=6pt,
}

\begin{document}

% ---- PORTADA ----
\thispagestyle{empty}
\begin{center}
  \vspace*{2cm}
  {\LARGE\bfseries Informe de Arquitectura Multiagente para \\[6pt]
  Teoría del Caso Aumentada}\\[2cm]
  {\large Santiago Esteban Redondo Perdomo \\[4pt]
  Juan David Conde Duarte}\\[1.5cm]
  {\large Ciencia de Datos}\\[1cm]
  {\large \today}\\[1cm]
  \vfill
\end{center}
\newpage

% ---- CUERPO ----

\section{Introducción}

El presente informe analiza la arquitectura de un sistema multiagente
diseñado para asistir en la construcción de la teoría del caso en el
ámbito del litigio jurídico. El sistema procesa documentos PDF de
demandas y escritos legales, extrayendo自动mente entidades
(actores, hechos, pruebas y normas), construyendo una matriz
Hecho-Prueba-Norma (HPN), generando una red multicapa de conocimiento,
calculando métricas estratégicas y simulando escenarios hipotéticos.
El frontend presenta un panel visual (dashboard) que traduce estos
datos en señales accionables para el litigante.

A continuación se presentan el análisis de utilidad, las limitaciones
identificadas, los errores encontrados durante el desarrollo y las
decisiones humanas que el sistema no puede reemplazar.

\section{Análisis de Utilidad}

\subsection{Automatización de la extracción de entidades}

El sistema emplea cuatro agentes especializados que reemplazan la
lectura manual del expediente:

\begin{itemize}[nosep]
  \item \textbf{Agente de Hechos} (\texttt{FactualAgent}): Divide el
    texto del PDF en oraciones, identifica fechas mediante
    expresiones regulares (\texttt{\textbackslash d\{1,2\}/\textbackslash
    d\{1,2\}/\textbackslash d\{4\}}) y filtra oraciones con más de
    30 caracteres como hechos jurídicos potenciales.
  \item \textbf{Agente de Actores} (\texttt{ActorsAgent}): Extrae
    nombres propios de personas usando el patrón
    \texttt{[A-ZÁÉÍÓÚÑ][a-záéíóúñ]+(?:\textbackslash s[A-ZÁÉÍÓÚÑ]
    [a-záéíóúñ]+)\{1,2\}}, filtrando términos legales (``código'',
    ``constitución'', ``juzgado'', etc.).
  \item \textbf{Agente de Pruebas} (\texttt{EvidenceAgent}): Detecta
    palabras clave probatorias (``contrato'', ``correo'',
    ``testimonio'', ``peritaje'', etc.) en el texto.
  \item \textbf{Agente Normativo} (\texttt{NormativeAgent}): Captura
    oraciones completas que contienen términos normativos (``ley'',
    ``artículo'', ``código'', ``constitución'') mediante el patrón
    \texttt{keyword + [\^{}.]\{0,120\}.}.
\end{itemize}

\subsection{Construcción de la matriz HPN}

El \texttt{HPNAgent} cruza hechos, pruebas y normas por índice
secuencial para construir la matriz Hecho-Prueba-Norma. Cada fila
recibe un estado epistémico (Completo si tiene prueba y norma;
Parcial en caso contrario) y un nivel de riesgo (Bajo si está
completo; Alto si es parcial). Esta matriz es la columna vertebral
del análisis y alimenta todos los componentes del dashboard.

\subsection{Generación de red multicapa}

El \texttt{NetworkAgent} construye un grafo dirigido con
\texttt{networkx} que organiza los nodos en cinco capas: actores,
hechos, pruebas, normas y riesgos. Las aristas representan relaciones
semánticas: \textit{soporta} (prueba \(\rightarrow\) hecho),
\textit{activa} (hecho \(\rightarrow\) norma), \textit{riesgo}
(riesgo \(\rightarrow\) hecho) y \textit{participa} (actor
\(\rightarrow\) hecho, solo si el nombre del actor aparece en el
texto del hecho). Esta estructura permite visualizar dependencias,
rutas críticas y el impacto de escenarios simulados.

\subsection{Métricas estratégicas}

El panel de métricas calcula seis indicadores compuestos a partir de
la matriz HPN y la estructura de red:

\begin{enumerate}[nosep]
  \item \textbf{Cobertura}: Porcentaje de filas con estado epistémico
    ``Completo''.
  \item \textbf{Fragilidad}: Porcentaje de filas con riesgo ``Alto''.
  \item \textbf{Centralidad}: Densidad relativa de la red
    (\(\frac{E}{N(N-1)}\)).
  \item \textbf{Contradicción}: Porcentaje de filas con
    contradicciones distintas de ``Ninguna''.
  \item \textbf{Robustez}: Complemento de la fragilidad (100\% --
    fragilidad).
  \item \textbf{Trazabilidad}: Porcentaje de filas con respaldo
    probatorio (prueba distinta de ``Sin prueba'').
\end{enumerate}

\subsection{Dashboard de señales accionables}

El frontend React traduce los datos en ocho pestañas funcionales:
Vista General (barra de preparación, señales accionables, entidades),
Matriz HPN (tabla interactiva con seis filtros), Red Multicapa (grafo
coloreado por capa con leyenda), Métricas (semáforo estratégico),
Alertas (cinco categorías: vacíos críticos, pruebas únicas,
contradicciones, baja trazabilidad, sin soporte probatorio),
Simulador (comparación antes/después con delta de nodos y aristas),
Preguntas (interrogatorio, contraexamen, cliente y perito) y Reporte
(exportación a PDF y HTML).

\subsection{Simulación de escenarios}

El \texttt{SimulatorAgent} ejecuta cuatro escenarios hipotéticos
sobre el grafo base: exclusión de prueba (elimina una arista),
nuevo peritaje (agrega nodo en capa pruebas), testigo contradictorio
(agrega nodo en capa pruebas) y nueva norma favorable (agrega nodo
en capa normas). Cada escenario reporta el cambio en número de nodos
y aristas, permitiendo al litigante evaluar el impacto antes de tomar
decisiones procesales.

\section{Limitaciones}

\subsection{Extracción basada exclusivamente en reglas}

Todos los agentes utilizan expresiones regulares y listas de palabras
clave. No emplean modelos de lenguaje (NLP profundo, transformers,
LLMs) ni aprendizaje automático. Esto limita la calidad de la
extracción:

\begin{itemize}[nosep]
  \item \texttt{FactualAgent} divide por puntos, lo que fragmenta
    incorrectamente oraciones con abreviaturas (``S.A.S.'',
    ``No.'') o numeración (``1.5'').
  \item \texttt{EvidenceAgent} solo detecta la primera ocurrencia de
    cada palabra clave por chunk, ignorando contextos donde la
    palabra clave aparece en un sentido no probatorio.
  \item \texttt{NormativeAgent} captura hasta 120 caracteres después
    de la keyword, pero no distingue entre una cita normativa
    genuina y una mención casual del término.
\end{itemize}

\subsection{Ausencia de un agente de contradicciones}

Aunque la matriz HPN incluye un campo \texttt{contradicciones}, este
siempre se inicializa como ``Ninguna''. No existe un agente dedicado
a detectar contradicciones entre hechos, entre pruebas o entre
declaraciones. El agente \texttt{auditor\_agent.py} existe como
archivo pero está vacío.

\subsection{Simulador superficial}

El \texttt{SimulatorAgent} modifica el grafo de forma elemental
(eliminar una arista, agregar un nodo) sin modelar el efecto
colateral que estos cambios tendrían en las métricas de red
(centralidad, densidad, PageRank). Tampoco considera la
jurisprudencia ni el contexto normativo para ponderar escenarios.

\subsection{Frontend sin estado persistente}

El dashboard React no persiste el estado del caso entre sesiones.
Cada carga de PDF comienza desde cero; no hay historial de casos,
comparación entre múltiples expedientes ni almacenamiento local.

\subsection{Dependencia de NetworkX para métricas}

El agente de métricas (\texttt{MetricsAgent}) depende de
\texttt{networkx} para betweenness centrality y PageRank. En grafos
muy pequeños o muy dispersos (como los generados a partir de un solo
documento), estas métricas pueden ser cero o no informativas, como
se observa en los datos de salida donde
\texttt{``betweenness''} contiene valores cercanos a cero y
\texttt{``pagerank''} es un diccionario vacío.

\section{Errores Encontrados y Corregidos}

\subsection{Clasificación incorrecta de actores}

\subsubsection*{Problema}

El patrón regex original del \texttt{ActorsAgent} capturaba cualquier
secuencia de dos o más palabras con inicial mayúscula, incluyendo
referencias legales como ``Código Civil'', ``La Constitución
Política'', ``Código General del Proceso'' y ``Corte Suprema''.
Estos términos no son personas jurídicas ni actores procesales, pero
se incluían en la lista de actores, inflando artificialmente el
conteo y distorsionando la red.

\subsubsection*{Solución}

Se implementó una lista de exclusión (\texttt{EXCLUDE\_WORDS}) con 35
términos legales comunes (``código'', ``constitución'', ``juzgado'',
``tribunal'', etc.) y se agregó una validación: si cualquiera de las
palabras del nombre candidato está en la lista de exclusión, se
descarta. Adicionalmente, se limitó el patrón a un máximo de tres
palabras (\texttt{\{1,2\}} en el cuantificador) para evitar capturar
frases completas.

\subsection{Actor único en la red multicapa}

\subsubsection*{Problema}

El \texttt{NetworkAgent} original conectaba únicamente al primer actor
de la lista (\texttt{actors[0]["nombre"]}) con todos los hechos,
ignorando por completo al resto de los actores. Esto producía un
grafo en estrella donde un solo nodo acumulaba todas las aristas y
los demás actores aparecían como nodos aislados.

\subsubsection*{Solución}

Se modificó el bucle de creación de aristas para iterar sobre todos
los actores. Para cada actor, se verifica si su nombre aparece en el
texto del hecho (comparación \texttt{lower()} de ambos lados). Solo
si existe coincidencia se crea la arista \textit{participa}. Esto
produce una red más realista donde cada actor se conecta únicamente
a los hechos que lo mencionan.

\subsection{Ausencia de normas en la respuesta del backend}

\subsubsection*{Problema}

El \texttt{NormativeAgent} extraía normas y las pasaba al
\texttt{HPNAgent}, pero el diccionario \texttt{entities} retornado
por la API no incluía la clave \texttt{"normas"}. El frontend
mostraba ``Normas Citadas'' vacío porque \texttt{entities.normas}
era un arreglo vacío.

\subsubsection*{Solución}

Se modificó \texttt{HPNService.generate()} para que, después de
extraer las normas, las agregue al diccionario \texttt{entities}
como \texttt{entities["normas"]} con la estructura
\texttt{\{"nombre": descripcion, "tipo": "norma"\}}. Dado que
\texttt{entities} es un diccionario mutable pasado por referencia, la
modificación persiste cuando se construye la respuesta JSON en
\texttt{app.py}.

\subsection{Normas genéricas (solo keyword)}

\subsubsection*{Problema}

El \texttt{NormativeAgent} original almacenaba únicamente la palabra
clave encontrada (``ley'', ``artículo'', ``código''), sin capturar
el contexto. En el dashboard aparecían normas como ``artículo'' en
lugar de ``artículo 1602 del Código Civil establece que todo
contrato...''.

\subsubsection*{Solución}

Se reemplazó la búsqueda por ocurrencia exacta con una expresión
regular que captura hasta 120 caracteres después de la keyword
(\texttt{re.escape(keyword) + r"[\^{}.]\{0,120\}\."}), extrayendo la
oración completa donde aparece la referencia normativa. Se agregó
un conjunto \texttt{seen} para evitar duplicados.

\subsection{Estilos Vite en conflicto con el dashboard}

\subsubsection*{Problema}

Los estilos CSS predeterminados de Vite (\texttt{index.css})
establecían un ancho fijo para \texttt{\#root}, centro de texto
obligatorio y variables de color que interferían con el layout del
dashboard.

\subsubsection*{Solución}

Se reescribió \texttt{index.css} eliminando todas las reglas
específicas de Vite y conservando solo un reseteo mínimo:
\texttt{box-sizing}, \texttt{font-family}, y
\texttt{min-height: 100vh} para \texttt{\#root}. Todos los estilos
del dashboard se definen exclusivamente en
\texttt{Dashboard.css}.

\section{Decisiones Humanas Requeridas}

\subsection{Validación de entidades extraídas}

El sistema no distingue entre un nombre de persona real y una entidad
mencionada incidentalmente. Por ejemplo, si el texto dice ``según la
Corte Suprema'', la palabra ``Corte'' hoy es filtrada, pero ``Suprema''
aislada podría no serlo. Un abogado debe revisar la lista de actores
extraídos y confirmar o descartar cada entrada.

\subsection{Clasificación del elemento jurídico}

La matriz HPN asigna automáticamente
\texttt{elemento\_juridico: "Pendiente Clasificación"} a todas las
filas. La decisión de clasificar cada hecho como pretensión,
excepción, fundamento o carga procesal es enteramente humana y debe
realizarse caso por caso.

\subsection{Evaluación de contradicciones}

El campo \texttt{contradicciones} se inicializa como ``Ninguna'' para
todas las filas. Detectar contradicciones reales entre hechos,
testimonios o pruebas requiere análisis jurídico contextual que el
sistema no puede realizar. El abogado debe identificar y documentar
cada contradicción manualmente.

\subsection{Determinación de acciones sugeridas}

La acción sugerida predeterminada (\texttt{"Validar soporte"}) es
genérica. El profesional debe evaluar cada fila de la matriz HPN y
determinar la acción concreta: interponer recurso, solicitar prueba
pericial, preparar contraexamen, etc.

\subsection{Interpretación de métricas estratégicas}

Las métricas (cobertura, fragilidad, robustez, etc.) son indicadores
cuantitativos que requieren interpretación jurídica cualitativa. Una
cobertura alta no significa necesariamente que el caso esté bien
fundamentado; puede indicar que las normas y pruebas asignadas son
genéricas. El abogado debe contextualizar cada métrica dentro de la
estrategia legal global.

\subsection{Selección y ponderación de escenarios}

El simulador ofrece cuatro escenarios predefinidos (exclusión de
prueba, nuevo peritaje, testigo contradictorio, nueva norma
favorable), pero la decisión de cuál simular primero, con qué
variantes y con qué peso estratégico depende del criterio del
litigante. El sistema no evalúa probabilidades ni recomienda
prioridades.

\subsection{Redacción de preguntas para audiencia}

El panel de preguntas sugeridas genera interrogantes automáticas
basadas en patrones de la matriz HPN. Sin embargo, la redacción
final, el orden de formulación y la adaptación al estilo del
juzgador son decisiones que solo el abogado puede tomar. Las
preguntas generadas son insumos, no sustitutos del cuestionario
definitivo.

\subsection{Exportación y confidencialidad}

El reporte exportable (PDF/HTML) contiene información sensible del
caso. La decisión de compartir, almacenar o imprimir el documento
debe cumplir con las normas de confidencialidad abogado-cliente y
protección de datos personales.

\section{Conclusiones}

La arquitectura multiagente implementada demuestra que es posible
automatizar etapas significativas del análisis jurídico preliminar:
extracción de entidades, construcción de matriz HPN, generación de
red multicapa, cálculo de métricas y simulación de escenarios. El
dashboard traduce estas estructuras de datos en señales accionables
que asisten al litigante en la preparación de su teoría del caso.

No obstante, el sistema presenta limitaciones importantes derivadas
de su enfoque basado exclusivamente en reglas: la extracción es
sensible al formato del documento de entrada, no detecta
contradicciones, no pondera escenarios con criterios jurídicos y
depende de validación humana en todos los puntos críticos. Los
errores encontrados (actores mal clasificados, red centrada en un
solo nodo, normas vacías o genéricas, estilos en conflicto) fueron
corregidos durante el desarrollo, pero evidencian la necesidad de
supervisión continua.

El sistema es una herramienta de aumentación, no de reemplazo. Su
valor reside en reducir el trabajo mecánico de extracción y
organización de información, permitiendo al abogado concentrarse en
las decisiones estratégicas de alto nivel que ningún algoritmo puede
suplir.

% ---- REFERENCIAS ----
\newpage
\begin{thebibliography}{99}

\bibitem{american2020}
American Psychological Association. (2020).
\textit{Publication manual of the American Psychological Association}
(7th ed.).

\bibitem{hagerty2023}
Hagerty, A. (2023). \textit{Multi-agent systems for legal reasoning:
A survey}. Artificial Intelligence and Law, 31(2), 245--278.
\url{https://doi.org/10.1007/s10506-022-09335-0}

\bibitem{networkx2014}
Hagberg, A., Schult, D., \& Swart, P. (2014).
\textit{Exploring network structure, dynamics, and function using
NetworkX}. Proceedings of the 7th Python in Science Conference
(SciPy 2008), 11--15.

\bibitem{pydantic2023}
Colvin, S. (2023). \textit{Pydantic: Data validation using Python
type hints} (Version 2.0). \url{https://docs.pydantic.dev/}

\bibitem{react2024}
Meta Platforms, Inc. (2024). \textit{React: A JavaScript library for
building user interfaces} (Version 19). \url{https://react.dev/}

\bibitem{vite2025}
Vite Team. (2025). \textit{Vite: Next generation frontend tooling}
(Version 8). \url{https://vite.dev/}

\bibitem{flask2024}
Pallets Project. (2024). \textit{Flask: Web development, one drop at
a time} (Version 3.0). \url{https://flask.palletsprojects.com/}

\end{thebibliography}

\end{document}
